%----------
%	ANALISIS GENERAL - LLMs
%----------	

The third and last classification approach that we followed was using Generative models to predict the "Análisis General" task. In particular, as seen in Table~\ref{table:classification_tasks}, a set of 64 experiments have been performed following the implementation in Section~\ref{implementation_genai}. A total of eight different windows have been defined as detailed in Table~\ref{tab:window-metrics-ageneral-genai}, corresponding to the two types of contexts and four prompts employed. Specifically, W1-W4 represented experiments using "Context\_0", while W5-W8 used "Context\_1." Within each context, the first window (W1 and W5) utilized "Prompt\_0\_EN", the second (W2 and W6) employed "Prompt\_0\_ES", the third (W3 and W7) applied "Prompt\_1\_EN", and the final window (W4 and W8) used "Prompt\_1\_ES".


\begin{table}[H]
    \centering
    \begingroup
    \small % Change to desired font size
    \ttabbox[\FBwidth]{
        \caption{AG Average Evaluation Metrics per Window in GenAI Experiments}
        \label{tab:window-metrics-ageneral-genai} % Unique label for the table
    }{
        \begin{tabularx}{\textwidth}{
            X
            X
            X
            >{\centering\arraybackslash}X
            >{\centering\arraybackslash}X
            >{\centering\arraybackslash}X
            >{\centering\arraybackslash}X
            >{\columncolor{gray!15}\centering\arraybackslash}X % Highlight this column
        }
            \hline
            \rowcolor{gray!30} % Adding a gray color to the header row
            \textbf{} & \textbf{Context} & \textbf{Prompt} & \textbf{Acc. (Std.)} & \textbf{Time (s)} & \textbf{W. Avg. F1-Score} & \textbf{C. Pos. F1-Score} & \textbf{C. Neg. F1-Score} \\
            \hline
            W1 & 0 & 0\_EN & 0.822 (0.171) & 630.238 & 0.814 & 0.880 & 0.650 \\
            \hline
            W2 & 0 & 0\_ES & 0.822 (0.190) & 680.053 & 0.811 & 0.877 & 0.678 \\
            \hline
            W3 & 0 & 1\_EN & 0.823 (0.177) & 641.584 & 0.815 & 0.880 & 0.662 \\
            \hline
            W4 & 0 & 1\_ES & 0.817 (0.201) & 688.112 & 0.804 & 0.872 & 0.680 \\
            \hline
            W5 & 1 & 0\_EN & 0.805 (0.214) & 647.818 & 0.790 & 0.861 & 0.672 \\
            \hline
            W6 & 1 & 0\_ES & 0.804 (0.218) & 685.949 & 0.790 & 0.860 & 0.680 \\
            \hline
            W7 & 1 & 1\_EN & 0.807 (0.244) & 657.664 & 0.792 & 0.862 & 0.677 \\
            \hline
            W8 & 1 & 1\_ES & 0.813 (0.211) & 700.921 & 0.800 & 0.867 & 0.686 \\
            \hline
            \multicolumn{8}{l}{Source: 'Training\_GenAI.xlsx'} \\
            \hline
        \end{tabularx}
    }
    \endgroup
\end{table}


Figure~\ref{fig:exp_analisis_general_classes_genai} presents the distribution of the experiments and their corresponding F1-Scores for the different aforementioned windows. The most notable insight that can be seen is that Context 0 (W1-W4) consistently achieved better F1 metrics compared to Context 1 (W1-W5). This is particularly interesting because, as defined in Section~\ref{contextual_information}, Context 0 contained no additional information, whereas Context 1 included detailed information about the \textit{"\#SeAcabó"} situation. Moreover, when analyzing the different windows within each context, the metrics for Context 0 tend to decrease across subsequent windows. However, performance for windows in Context 1 showed an improving trend. This suggests that the absence of context led to better overall results for this task, as both accuracy (0.823 in W3) and weighted average F1-Scores (0.880 in W1 and W3) were higher for each window in Context 0 compared to the best-performing window in Context 1. Additionally, prompts in English yielded better weighted average and "Comentario Positivo" F1-Scores in Context 0, while Spanish prompts performed better in terms of "Comentario Negativo" F1-Score in Context 1.

Moreover, we can also observe the distribution of experiments and metrics for both classes. The models consistently performed better on the "Comentario Positivo" class with F1-Scores around 0.87 compared to 0.67 for "Comentario Negativo". However, there was a positive trend in the "Comentario Negativo" class, with F1 scores improving from 0.680 in W4 of Context 0 to 0.686 in W8 of Context 1.


\begin{figure}[H]
	\ffigbox[\FBwidth]
	{\caption{Distribution of the GenAI Experiments over Time for "Análisis General", per Class with the Different Windows, and Trend Lines for Each Class}\label{fig:exp_analisis_general_classes_genai}}
 % [Aquí ponemos como queremos que aparezca en el índice, sin la referencia]{Aquí como queremos que aparezca en el documento, con la referencia}
	{\includegraphics[scale=0.25]{Plantilla_TFG_ingles_2019/imagenes/GenAI_Experiments in _Análisis General_ f1score.pdf}}
\end{figure}

\newpage

%\begin{figure}[H]
%	\ffigbox[\FBwidth]
%	{\caption{Distribution of the GenAI Experiments over Time for "Análisis General" with the Different Windows and Weighted Avg. F1-Score Trend Line}\label{fig:exp_analisis_general_genai}}
 % [Aquí ponemos como queremos que aparezca en el índice, sin la referencia]{Aquí como queremos que aparezca en el documento, con la referencia}
%	{\includegraphics[scale=0.25]{Plantilla_TFG_ingles_2019/imagenes/GenAI_Experiments in _Análisis General_ weighted.pdf}}
%\end{figure}





Table~\ref{tab:ageneral-genai-bestmodels_metrics} provides a summary of the top 5 best-performing generative models for the "Análisis General" task. Overall, we observe that the accuracy and F1-Score were consistent around 0.81 and 0.80 respectively, with the standard deviation reaching approximately 0.21, which was higher than that of other classification approaches described in Table~\ref{tab:ageneral-bestmodels_metrics} for ML and Table~\ref{tab:ageneral-dl-bestmodels_metrics} for DL.


\begin{table}[H]
    \centering
    \begingroup
    \small % Change to desired font size
    \ttabbox[\FBwidth]{
        \caption{AG Evaluation Metrics for the Best GenAI Models}
        \label{tab:ageneral-genai-bestmodels_metrics} % Unique label for the table
    }{
        \begin{tabularx}{\textwidth}{
            X
            >{\centering\arraybackslash}X
            >{\centering\arraybackslash}X
            >{\centering\arraybackslash}X
            >{\columncolor{gray!15}\centering\arraybackslash}X % Highlight this column
        }
            \hline
            \rowcolor{gray!30} % Adding a gray color to the header row
            \textbf{} & \textbf{Acc. (Std.)} & \textbf{Weighted Avg. F1-Score} & \textbf{C. Positivo F1-Score} & \textbf{C. Negativo F1-Score} \\
            \hline
            \multicolumn{5}{l}{\textbf{Naive Models}} \\
            \hline
            \agmajor & 0.778 (0.010) & 0.680 & 0.875  & 0.000 \\
            \hline
            \agrandom & 0.654 (0.012) & 0.654 & 0.778 & 0.222 \\
            \hline
            \multicolumn{5}{l}{\textbf{Zero-Shot Models}} \\
            \hline
            \agmistralsevenb & 0.815 (0.211) & 0.802 & 0.868 & 0.690 \\
            \hline
            \agllamathreeoneone & 0.815 (0.211) & 0.801 & 0.868 & 0.689 \\
            \hline
            \aggemmatwob & 0.814 (0.212) & 0.800 & 0.868 & 0.688 \\
            \hline
            \agllamathree & 0.815 (0.209) & 0.843 & 0.869 & 0.687 \\
            \hline
            \agllamathreeonefive & 0.813 (0.212) & 0.800 & 0.867 & 0.687  \\
            \hline

            \multicolumn{5}{l}{Source: 'Training\_GenAI.xlsx'} \\
            \hline
        \end{tabularx}
    }
    \endgroup
\end{table}


Based on this analysis, model \textbf{\agmistralsevenb} was selected as the best-performing model for this task and approach. This decision was chosen because of its high F1-score for the "Comentario Negativo" class, which was the highest among the top-performing models. In Figure~\ref{fig:exp_ageneral_genai_bestmodel} we can see in more detail the metrics of this model. In particular, the model correctly classified 1006 instances of "Comentario Positivo" with only 28 misclassifications, while it struggled more with the "Comentario Negativo" class, correctly classifying 339 instances but misclassifying 277 instances. In this case, these results cannot be attributed to "Comentario Positivo" being the majority class, because as explained in Section~\ref{zero_shot}, the zero-shot learning approach means the model had no prior exposure to training examples. Therefore, while the model outperformed the two naive baseline models and demonstrated solid performance in classifying positive comments, identifying negative comments remains a more challenging task for the zero-shot capabilities of this type of generative models. 


% Best Model
\begin{figure}[H]
	\ffigbox[\FBwidth]
	{\caption{Confusion Matrix for the Best GenAI Model (\textbf{\agmistralsevenb}) for "Análisis General"}\label{fig:exp_ageneral_genai_bestmodel}}
 % [Aquí ponemos como queremos que aparezca en el índice, sin la referencia]{Aquí como queremos que aparezca en el documento, con la referencia}
	{\includegraphics[scale=0.15]{Plantilla_TFG_ingles_2019/imagenes/GenAI_Experiments in _Análisis General_ best_model.pdf}}
\end{figure}

\newpage